\section{A cortico-basal-ganglia-collicular loop}

The controlled comparison isolates two structural conditions that are jointly necessary for learning: the policy and value computations must be coupled through a shared, jointly-maintained recurrent memory (V2 fails without it), and the grounded \emph{priority} stream must command the actor (V3 fails when the \emph{value} stream does instead). Neither condition is an arbitrary engineering choice. Each maps onto a well-characterised feature of the cortico-basal-ganglia-collicular machinery that the primate brain uses to convert sensory evidence into committed, value-weighted action. We develop that mapping here, both because it motivates the architecture and because it makes the two failure modes biologically interpretable rather than incidental.

\subsection{Actor and critic as a striatal division of labour}

The model factors the decision into an actor that emits the wait/declare policy and a distributional critic that evaluates the state. This factorisation has a direct anatomical precedent. In the basal ganglia, the dorsal striatum is most closely associated with the actor---the mapping from state to action---while the ventral striatum, together with its dopaminergic input, is most closely associated with the critic, the evaluation of state value and the computation of a prediction error \citep{odoherty2004dissociable,maia2009reinforcement}. Human imaging during instrumental learning dissociates these roles: ventral striatal activity tracks reward prediction in both Pavlovian and instrumental settings, whereas dorsal striatal activity is selective for the instrumental, action-contingent case \citep{odoherty2004dissociable}. The two systems are anatomically distinct but functionally interdependent, which is exactly the relationship the model instantiates between its separate actor and critic heads sitting on one shared working memory.

We stress that this correspondence is an organising idealisation rather than settled anatomy. A literal localisation of ``the critic'' to a single nucleus overstates the cleanliness of the division: action-value and reward signals are distributed broadly across the striatum, and the dorsal/ventral split is better read as a gradient than a partition \citep{maia2009reinforcement}. We therefore use the actor-critic mapping as a conceptual frame for the model's labour division, not as a claim that the priority stream $\mu$ and the value stream $q$ are two anatomically separable attention nodes.

\subsection{Dopamine as the coupling that makes V2's collapse sensible}

The sharpest empirical result is that severing the coupling between the policy and value computations abolishes learning entirely: the independent variant (V2), in which the actor acts from a private memory and the critic evaluates a separate private memory, collapses to an always-wait policy with a hit rate of $0.000$, never once declaring a change across $500$ fresh episodes. This is not a capacity failure---the network retains the parameters and pathways to express the task---but a credit-assignment failure. When the actor and critic occupy different representations, the advantage handed to the actor is computed on a state the actor does not occupy, credit is systematically misattributed, and the policy settles into the only basin that is safe under any value estimate: never act, and never be penalised.

The biological counterpart of the required coupling is the dopaminergic teaching signal. Midbrain dopamine neurons in the substantia nigra pars compacta (SNc) and the ventral tegmental area encode a temporal-difference reward-prediction error: they respond to unexpected reward, shift their response to the earliest reward-predicting cue as learning proceeds, and are suppressed when a predicted reward is omitted \citep{schultz1997neural,hollerman1998dopamine}. This error is precisely the quantity the critic computes, and it is broadcast to striatal synapses where it gates plasticity. The corticostriatal learning rule is a three-factor rule---presynaptic activity, postsynaptic activity, and dopamine concentration jointly determine whether a corticostriatal synapse potentiates or depresses---so that the policy encoded in striatal weights is updated only in proportion to the dopaminergic error supplied by the value system \citep{morita2014striatal}. Value and policy are thereby bound into a single closed loop: the actor learns nothing except through the critic's broadcast error.

This is the conceptual precedent for V2. An agent whose policy circuit receives no teaching signal derived from the same evaluated state has no instructive gradient with which to discover that declaring at change onset is rewarded; the always-wait collapse is the rational endpoint of an actor that cannot tell when acting helps. In the model, the shared deep memory $H_2$ plays the role the dopaminergic loop plays in the brain: because the actor's own output writes $H_2$ and the critic reads the same $H_2$, ``pressing now looks good'' is a statement about the state the actor actually occupies, and credit lands where it was earned. Removing the shared memory is the in-silico analogue of opening the dopaminergic teaching loop.

The parallel must be drawn with care, and we flag the principal disanalogy explicitly. In the brain the critic-to-actor influence is a \emph{broadcast scalar}---a diffusely released dopamine prediction error flowing from the critic's evaluation back to the actor's synapses. In the model the coupling is a \emph{shared recurrent state} that the actor writes and both streams read. The substrate and the dominant direction of the brain's coupling (critic $\rightarrow$ actor, via dopamine) differ from the model's most visible coupling (priority $\rightarrow H_2$, a written state). We present the correspondence as motivation, not isomorphism. We note, however, that both directions genuinely co-exist in the model: the priority stream's write into $H_2$ closes the loop forward, while the critic's gradient---the model's analogue of the dopamine error---closes it backward by teaching the actor. It is, in both systems, a loop; what we decline to claim is that the loop is realised by the same machinery.

We further note that some of the strongest evidence that value and decision signals can share a substrate comes from parietal cortex: neurons in the lateral intraparietal area (LIP) carry both the accumulating decision variable and reward-magnitude or reward-probability signals on the same population, with reward entering as an offset to the decision process \citep{platt1999neural,sugrue2004matching}. That a single neural substrate intermixes value and decision lends biological plausibility to the model's premise that one shared working memory can support both an actor and a critic read-out.

\subsection{The superior colliculus as the grounded node that must command action}

The second necessary condition---revealed by the swapped variant (V3), which is architecturally identical to the working model but routes the value stream to the actor and the priority stream to the critic, and which likewise collapses to a hit rate of $0.000$---is that the actor must be driven by the image-grounded stream. In V3 the loop is closed and credit assignment is clean, yet the actor is fed a representation whose residual is memory rather than the live image; it carries almost no current visual content, so there is nothing in the actor's input from which to detect the change. The grounded stream is what the policy must read.

This condition maps onto the role of the superior colliculus (SC), and of the visual cells of the frontal eye field (FEF), as the grounded priority node that commits the response. The SC is causally necessary for covert spatial attention: reversible inactivation of the intermediate SC produces a marked deficit in covert target selection even when the eyes never move, and does so without changing the attentional gain measured in extrastriate cortex \citep{zenon2012attention,lovejoy2010inactivation}. The colliculus is best understood not as a gain knob on sensory cortex but as the threshold that converts accumulating cortical evidence into a committed choice: collicular activity reflects the commitment to a response rather than the evidence accumulation itself \citep{stine2023differential}. Its causal weight is heavy---SC inactivation is more devastating to covert attentional performance than FEF inactivation \citep{bollimunta2018fefsc}---and its influence extends into the basal ganglia: attentional modulation in the caudate nucleus depends on an intact SC, defining an SC-caudate-cortex loop in which the collicular priority signal conditions striatal processing \citep{herman2020attention}. That loop is a loose biological analogue of the model's arrangement, in which the priority stream both commits the action and writes the shared state that the value circuit then reads.

The complementary cortical evidence points the same way. During covert attention the FEF carries the attentional signal on its \emph{visual} (image-grounded) cells, while its movement-related cells are suppressed---the grounded representation, not the motor command, does the perceptual work \citep{thompson2005neuronal}. Subthreshold microstimulation of the FEF, below the level that evokes a saccade, enhances visual sensitivity at the corresponding location, establishing the FEF as a controller that acts through the grounded visual pathway \citep{moore2001control,moore2004microstimulation}. More generally, perceptual decisions are formed by integrating grounded sensory evidence to a bound \citep{gold2007neural}, and the grounded representation is the indispensable input: inactivating sensory cortex (MT) impairs the discrimination itself, whereas inactivating the downstream accumulator (LIP) leaves the discrimination largely intact \citep{katz2016dissociated}. This dissociation is the single closest biological parallel to V3---an actor cut off from the grounded sensory representation cannot perform the discrimination no matter how well-formed its downstream evaluation.

We again temper the mapping. The model's priority stream $\mu$ corresponds most naturally to the FEF-visual and SC commitment integrator, and the value stream $q$ to the value/critic system in the basal ganglia and its dopaminergic input; the latter is not itself an attention node. SC, FEF, and LIP are each single priority maps, not actor-critic pairs, and we do not claim that $\mu$ and $q$ map onto two distinct attention areas. The claim is narrower and, we think, robust: a grounded, collicular-style commitment node must be the one that drives the action, and a value circuit must be the one that evaluates it.

\subsection{The loop, stated whole}

The three results assemble into a single closed loop. Grounded priority---visual cortex feeding the SC and the FEF-visual cells---commits the declare response and writes the shared working state. The value system in the striatum and basal ganglia reads that state and evaluates it. The evaluation is broadcast back as an SNc dopamine temporal-difference error that teaches the actor's plasticity, closing the loop. Only when the loop is closed (a shared state coupling policy and value) \emph{and} the grounded node commands the action does the system learn. V2 corresponds to an open loop: the policy circuit receives no teaching signal tied to the evaluated state, as if the dopaminergic coupling reaching the actor were severed. V3 corresponds to a loop in which the value circuit wrongly issues the motor commitment without grounded collicular evidence---credit assignment is intact, but the commanding representation cannot see the world. The working model is the one configuration that respects both the closure and the grounding of the biological circuit, and it is the only one of the three that learns the task.
